% ── Section: Implementation and Real-Time Deployment ─────────────────
% NEW 27/8/2026, written for the journal venue. A conference version had no
% room for this; the journal does, and for JRTIP it is the section that
% decides whether the paper belongs there at all.
% Every number here comes from experiments/realtime/evidence.md.
% Open: the accuracy axis of Table~\ref{tab:pareto} is not yet measured --
% it needs a detection-head model trained on the same data (requested) and
% the dense GT re-review (in progress). Do not submit with that gap.
% ─────────────────────────────────────────────────────────────────────

\section{Implementation and Real-Time Deployment}
\label{sec:implementation}

\subsection{Computational Cost of the Proposed Components}
\label{subsec:complexity}

The CTRV filter of Sec.~\ref{subsec:c2} replaces a linear Kalman filter on a
$5$-dimensional state, so its per-track cost is dominated by the covariance
propagation $\mathbf{F}\mathbf{P}\mathbf{F}^{\top}$ and the Joseph-form
update: $\mathcal{O}(d^{3})$ with $d=5$, against $d=8$ for the
constant-velocity filter it replaces on motorcycles.
The Jacobian (Eq.~\ref{eq:ctrv_x}--\ref{eq:ctrv_theta}) is evaluated in closed
form, adding no iteration.
Total tracking cost is therefore $\mathcal{O}(n)$ in the number of tracks for
prediction, plus the $\mathcal{O}(nm)$ association already present in
ByteTrack for $n$ tracks and $m$ detections; the uncertainty gate adds one
$2\times2$ Mahalanobis test per lost motorcycle track, and NSA a single scalar
scaling of $\mathbf{R}$ per detection.
The context estimator (C1) is four scalar statistics over the current frame
and track set, and the reliability score (C3) is a weighted sum of five
per-track quantities already maintained by the tracker; neither introduces a
term that grows faster than $\mathcal{O}(n)$.

Memory is likewise bounded: the tracker holds one $5\times5$ covariance per
motorcycle track and one $8\times8$ per other vehicle, so at the $\sim$47
tracks/frame of the \textsc{dense} window the filter state is under
$20$\,kB in total.

\subsection{End-to-End Latency}
\label{subsec:endtoend}

Sec.~\ref{subsec:runtime} reported tracker-side latency with the detector
excluded, which isolates what the proposed modules cost but says nothing about
throughput.
Table~\ref{tab:endtoend} reports the whole path --- decode, context,
detection, tracking, reliability scoring --- on the same CPU, one variant per
process, warm-up frames discarded.

\begin{table}[h]
  \centering
  \caption{End-to-end per-frame cost, detector included
           (100 frames per row, one process per row, Intel Core i5-12500H,
           CPU only). Source is $25$\,fps, so the real-time budget is
           $40$\,ms.}
  \label{tab:endtoend}
  \begin{tabular}{lccccc}
    \toprule
    Configuration & detect & track & total & p95 & fps \\
    \midrule
    Exp~0 (baseline)     & 109.8 & 5.8 & 117.1 & 131.5 & 8.5 \\
    Exp~1 (+C2)          & 109.5 & 9.5 & 120.4 & 135.0 & 8.3 \\
    Exp~2 (+C1, C2)      & 110.9 & 7.8 & 120.1 & 140.8 & 8.3 \\
    Exp~3 (+C1, C2, C3)  & 110.4 & 7.4 & 120.3 & 133.3 & 8.3 \\
    \bottomrule
  \end{tabular}
\end{table}

Two things change once the detector is inside the measurement.

\paragraph{The proposed components are nearly free, and the detector is not}
C2 costs $+22\%$ of tracker time (Table~\ref{tab:runtime}) but $+2.8\%$ of a
frame --- $3.3$\,ms on $117$\,ms.
C1 and C3 fall inside run-to-run noise.
Detection accounts for $91\%$ of the frame in the shipped configuration.
Any discussion of real-time feasibility for this pipeline is a discussion
about the detector's operating point; the tracking contribution of this paper
is not the term that decides it.

\paragraph{The shipped configuration is not real-time, for an avoidable reason}
At $8.3$\,fps the proposed configuration misses the $40$\,ms budget by
$3\times$.
The cause is not the detector's scale but its \emph{head}: the deployed
checkpoint is a YOLO11n \emph{segmentation} model, while the pipeline consumes
only bounding boxes.
A detection head of the same scale costs $58.2$\,ms against $130.7$\,ms at
$640$\,px --- a factor of $2.25$ spent computing a $160\times160\times32$
prototype tensor that is discarded on every frame.

\subsection{Reaching the Real-Time Budget}
\label{subsec:pareto}

Table~\ref{tab:pareto} sweeps the detector operating point with a detection
head.
The budget is met at $448$\,px and below.

\begin{table}[h]
  \centering
  \caption{Detector operating points (detection head, PyTorch, CPU).
           Accuracy column pending --- see the caveat below.}
  \label{tab:pareto}
  \begin{tabular}{lcccc}
    \toprule
    Input size & detect (ms) & total (ms) & p95 & fps \\
    \midrule
    $640$ & 58.2 & 62.9 & 69.6 & 15.9 \\
    $448$ & 24.3 & 26.6 & \textbf{32.9} & \textbf{37.6} \\
    $320$ & 17.3 & 19.5 & 24.7 & 51.4 \\
    $256$ & 14.2 & 16.3 & 20.5 & 61.4 \\
    \bottomrule
  \end{tabular}
\end{table}

\todo{Accuracy axis. Table~\ref{tab:pareto} is half a Pareto curve: it has
cost and no quality. Filling it needs (i) a detection-head model trained on
the same data as the deployed segmentation checkpoint, and (ii) the
\textsc{dense} ground-truth re-review, since no accuracy figure should be
quoted before that lands. Do not submit without this.}

\subsection{Export Backends}
\label{subsec:backends}

The usual route to CPU real-time is to export the network to an inference
runtime.
On this machine it is a regression in every combination we measured:
OpenVINO costs $+25\%$ over PyTorch for both the segmentation head
($163.9$ vs.\ $130.7$\,ms) and the detection head ($73.0$ vs.\ $58.2$\,ms),
and ONNX Runtime $+75\%$ ($228.7$\,ms).
Half precision is slower still ($79.3$\,ms), as expected on a CPU with no
native half-precision path, where the conversion is real work and the
arithmetic is not cheaper.

The consistency of the OpenVINO penalty across both heads suggests the
backend rather than one model, and the likely mechanism is that only the
backbone traverses the exported graph while postprocessing remains in the
framework --- giving up fused CPU kernels without replacing the dominant term.
We report this as measured, with default export settings and no thread or
stream tuning, rather than as a general claim about these runtimes.
Integer quantization was not evaluated: it requires a calibration set from the
detection training data, which we do not hold.

The practical conclusion for deployment on commodity CPUs is the reverse of
the usual advice: choosing the right \emph{head} was worth $2.25\times$, while
changing the \emph{runtime} cost $25$--$75\%$.

\subsection{Measurement Protocol}
\label{subsec:protocol}

Two protocol choices are worth stating because violating either produces
differences larger than the ones under study.

First, each configuration is timed in a separate process.
Timing several in one run makes the measurement track position rather than
configuration: the same baseline code measured $13.19$\,ms first and
$4.07$\,ms last, a $3.2\times$ spread that exceeds every difference in
Table~\ref{tab:endtoend} and flatters whichever configuration runs last.

Second, warm-up frames are discarded rather than averaged in; the first frames
pay for lazy initialization and kernel autotuning, and we measure a mean of
$219$\,ms across them against $143$\,ms in steady state.

We also caution that absolute end-to-end numbers are comparable within a sweep
but not across sessions: the same detector on the same footage measured
$109.8$\,ms in one sweep and $130.7$\,ms in another.
Every comparison drawn above is within a single sweep.
